\subsubsection{One-Point Test for Anomaly Classes and Coordinate Representation of Cross-Anomalies: Torsion-Free Criteria and Euler Quotient Invariants}\label{subsubsec:pom-anom-torsionfree-euler-quotient}
This paragraph supplements two ``constant-level'' corollaries on top of the anomaly cohomology framework for the three-generator fragment (Corollary \ref{cor:pom-anom-cohomology-class}) and the moment amplification law (Proposition \ref{prop:pom-anom-moment-amplification}):
the first provides a \emph{one-point zero-anomaly test} in the torsion-free case;
the second organizes the higher-order cross-anomaly family \(\{\Delta_S\}\) into \emph{minimal linear coordinates} for Euler/direct-product decomposability (quotient space isomorphism and dimension formula).

\begin{corollary}[One-point zero-anomaly test in the torsion-free case]\label{cor:pom-anom-one-point-zero-test}
Let \(H\) denote the anomaly cohomology group of the projection word fragment (the abelian group obtained by quotienting the central label group by coboundaries), and let \([\mathrm{Anom}(w)]\in H\) be the anomaly class of any \(1\)-morphism \(w\).
If \(H\) is torsion-free (i.e., \(qa=0\Rightarrow a=0\) for all integers \(q\ge 1\)), then for any integer \(q\ge 2\) the following strict equivalence holds:
\[
\boxed{\;
[\mathrm{Anom}(w)]=0\quad\Longleftrightarrow\quad [\mathrm{Anom}(\mathrm{MOM}_q\circ w)]=0\; }.
\]
Therefore, on such groups \(H\), the ``zero anomaly'' condition can be completely determined by verification at any single moment order.
\end{corollary}
\begin{proof}
By Proposition \ref{prop:pom-anom-moment-amplification},
\(
[\mathrm{Anom}(\mathrm{MOM}_q\circ w)]=q\,[\mathrm{Anom}(w)]
\).
If \([\mathrm{Anom}(w)]=0\) then the right-hand side is zero, hence \([\mathrm{Anom}(\mathrm{MOM}_q\circ w)]=0\).
Conversely, if \([\mathrm{Anom}(\mathrm{MOM}_q\circ w)]=0\), then \(q\,[\mathrm{Anom}(w)]=0\), and the torsion-free property yields \([\mathrm{Anom}(w)]=0\). \qed
\end{proof}

\begin{theorem}[Coprime Moment Order Zero Vanishing Principle: Two-Point Completeness of Anomaly Verification (Without Torsion-Free Assumption)]\label{thm:pom-anom-coprime-two-point-test}
Suppose that the anomaly class under the PW semantics lies in some abelian group \(H\) and satisfies the moment amplification law: for any integer \(q\ge 2\),
\[
[\mathrm{Anom}(\mathrm{MOM}_q\circ w)]=q\cdot[\mathrm{Anom}(w)]\in H
\]
(Proposition \ref{prop:pom-anom-moment-amplification}).
If there exist coprime integers \(q_1,q_2\ge 2\) such that
\[
[\mathrm{Anom}(\mathrm{MOM}_{q_1}\circ w)]=0,\qquad
[\mathrm{Anom}(\mathrm{MOM}_{q_2}\circ w)]=0,
\]
then necessarily \([\mathrm{Anom}(w)]=0\).
Consequently, for all \(q\ge 2\) we have \([\mathrm{Anom}(\mathrm{MOM}_{q}\circ w)]=0\).
\end{theorem}
\begin{proof}
Let \(a:=[\mathrm{Anom}(w)]\in H\). By the amplification law, \(q_1 a=0\) and \(q_2 a=0\).
Take integers \(u,v\) such that \(u q_1+v q_2=1\) (B\'ezout's identity), then
\[
a=(u q_1+v q_2)a=u(q_1 a)+v(q_2 a)=0.
\]
Substituting back into the amplification law yields the conclusion for all \(q\ge 2\). \qed
\end{proof}

\begin{theorem}[Anomaly--Oracle Collapse: Infinite Moment Order Audit \texorpdfstring{$\Longleftrightarrow$}{<->} Two Coprime Moment Order Audits]\label{thm:pom-anom-all-moments-collapse-to-two}
Let the proposition
\[
\mathcal P(w):\quad \forall q\ge 2,\;[\mathrm{Anom}(\mathrm{MOM}_q\circ w)]=0
\]
denote ``anomaly vanishes at all moment orders.''
Then \(\mathcal P(w)\) is equivalent to the existence of coprime integers \(q_1,q_2\ge 2\) such that
\(
[\mathrm{Anom}(\mathrm{MOM}_{q_1}\circ w)]=[\mathrm{Anom}(\mathrm{MOM}_{q_2}\circ w)]=0
\).
\end{theorem}
\begin{proof}
``\(\Rightarrow\)'' is obvious (e.g., take \((q_1,q_2)=(2,3)\)).
``\(\Leftarrow\)'' follows immediately from Theorem \ref{thm:pom-anom-coprime-two-point-test}. \qed
\end{proof}

\begin{corollary}[Auditable Trigger Condition for Oracle Collapse: Two-Point Finite Certificate for All-Moment Zero Anomaly]\label{cor:pom-anom-oracle-collapse-two-point-trigger}
Any condition that relies logically on proposition \(\mathcal P(w)\) as a sufficient condition—thereby depending on ``infinite audit over all moment orders \(q\ge 2\)''—can be compressed to verifying only two coprime moment orders; in other words, the audit complexity of the ``all-moment anomaly axis'' collapses logically to two points.
\end{corollary}

\begin{theorem}[Torsion Eliminability Test and Minimal Elimination Order]\label{thm:pom-anom-torsion-eliminability-min-order}
Under the setting of Theorem \ref{thm:pom-anom-coprime-two-point-test}, let \(a:=[\mathrm{Anom}(w)]\in H\), and denote
\[
H_{\mathrm{tor}}:=\{x\in H:\ \exists\,n\ge 1\ \text{such that}\ nx=0\}
\]
as the torsion subgroup of \(H\).
Then:
\begin{enumerate}
  \item There exists an integer \(q\ge 2\) such that \([\mathrm{Anom}(\mathrm{MOM}_q\circ w)]=0\) if and only if \(a\in H_{\mathrm{tor}}\).
  \item If \(a\notin H_{\mathrm{tor}}\) and \(H\) can be embedded as an additive subgroup of some normed real linear space \((V,\|\cdot\|)\), then
  \[
  \bigl\|[\mathrm{Anom}(\mathrm{MOM}_q\circ w)]\bigr\|
  =\|qa\|\xrightarrow[q\to\infty]{}\infty.
  \]
  \item If \(a\in H_{\mathrm{tor}}\) with order \(\mathrm{ord}(a)\in\NN\), then for any integer \(q\ge 2\) the following strict equivalence holds:
  \[
  \boxed{\;
  [\mathrm{Anom}(\mathrm{MOM}_q\circ w)]=0\quad\Longleftrightarrow\quad \mathrm{ord}(a)\mid q\; }.
  \]
  In particular, \(\mathrm{ord}(a)\) is the minimal positive integer (minimal elimination order) for which the anomaly can be eliminated by moment compilation.
\end{enumerate}
\end{theorem}
\begin{proof}
By Proposition \ref{prop:pom-anom-moment-amplification},
\(
[\mathrm{Anom}(\mathrm{MOM}_q\circ w)]=qa
\).
Hence there exists \(q\ge 2\) making this zero if and only if there exists \(q\ge 2\) such that \(qa=0\), which is equivalent to \(a\in H_{\mathrm{tor}}\).
If \(a\notin H_{\mathrm{tor}}\) and \(H\hookrightarrow (V,\|\cdot\|)\), then \(a\neq 0\) and \(\|qa\|=q\|a\|\to\infty\).
If \(a\in H_{\mathrm{tor}}\) with \(\mathrm{ord}(a)=n\), then \(qa=0\) if and only if \(n\mid q\), so the minimal positive integer is precisely \(n\). \qed
\end{proof}

\begin{corollary}[Finite Covering Elimination and Least Common Multiple Law for Torus Phase Anomalies]\label{cor:pom-anom-torus-phase-elimination}
Suppose the anomaly takes values in the compact torus
\(
H=(\RR/2\pi\ZZ)^r
\), and write \(a=(a_1,\dots,a_r)\in H\).
Then there exists an integer \(q\ge 2\) such that \(qa=0\) in \(H\) if and only if for each \(1\le j\le r\),
\[
a_j\in 2\pi\QQ/2\pi\ZZ.
\]
In this case, if each \(j\) admits the reduced representation
\(
a_j\equiv 2\pi\,u_j/v_j\pmod{2\pi}
\) (with \(\gcd(u_j,v_j)=1\)), then the minimal elimination order is
\[
\boxed{\;
q_\star=\mathrm{lcm}(v_1,\dots,v_r)\; },
\]
and for any integer \(q\ge 1\), \(qa=0\) if and only if \(q_\star\mid q\).
\end{corollary}
\begin{proof}
In \(H\), \(qa=0\) if and only if for each \(j\), \(q a_j\equiv 0\ (\mathrm{mod}\ 2\pi)\), i.e., \(a_j/(2\pi)\in \QQ/\ZZ\) and its denominator divides \(q\).
The minimal positive integer is given by the least common multiple of the denominators. \qed
\end{proof}

\begin{corollary}[Torsion Sufficient Condition for Roots-of-Unity Phases]\label{cor:pom-anom-cyclotomic-phase-torsion}
Under the setting of Corollary \ref{cor:pom-anom-torus-phase-elimination}, if there exists a positive integer \(N\) such that for each coordinate \(a_j\), \(N a_j\equiv 0\ (\mathrm{mod}\ 2\pi)\) (equivalently, \(e^{\mathrm{i}a_j}\) is a finite-order root of unity), then \(a\in H_{\mathrm{tor}}\), so there exists a finite order \(q\ge 2\) such that \([\mathrm{Anom}(\mathrm{MOM}_q\circ w)]=0\).
\end{corollary}
\begin{proof}
By hypothesis, one can take the same \(N\) such that \(Na=0\) in \(H\), hence \(a\) is a torsion element.
The conclusion then follows from Theorem \ref{thm:pom-anom-torsion-eliminability-min-order}. \qed
\end{proof}

\begin{corollary}[Unbounded Amplification and Threshold Crossing for Non-Torsion Anomalies]\label{cor:pom-anom-unbounded-amplification}
Under the setting of Theorem \ref{thm:pom-anom-torsion-eliminability-min-order}, if the anomaly class \(a=[\mathrm{Anom}(w)]\) is not a torsion element, then the set
\(
\{[\mathrm{Anom}(\mathrm{MOM}_q\circ w)]:q\ge 2\}
\)
is infinite in \(H\).
Furthermore, if \(H\) can be viewed as an additive subgroup of some normed real linear space \((V,\|\cdot\|)\) (for example, if \(H\) itself is a real vector space or an additive subgroup thereof), with the norm restricted to \(|\cdot|\), then for any threshold \(\tau>0\) there exists an integer \(q\ge 2\) such that
\[
\boxed{\;
\bigl|[\mathrm{Anom}(\mathrm{MOM}_q\circ w)]\bigr|>\tau\; }.
\]
\end{corollary}
\begin{proof}
By Proposition \ref{prop:pom-anom-moment-amplification}, for \(q\ge 2\) we have
\(
[\mathrm{Anom}(\mathrm{MOM}_q\circ w)]=qa
\).
If there exist \(q_1\neq q_2\) such that \(q_1a=q_2a\), then \((q_1-q_2)a=0\), forcing \(a\in H_{\mathrm{tor}}\), contradicting the hypothesis; hence \(\{qa:q\ge 2\}\) is infinite.
Under the normed embedding, homogeneity gives \(|qa|=q|a|\to\infty\), so one may take \(q>\tau/|a|\) to achieve threshold crossing. \qed
\end{proof}

\begin{proposition}[Support constraint on cross-anomalies: vanishing at trivial coordinates]\label{prop:pom-cross-anom-support-constraint}
Under the setting of Definition \ref{def:pom-higher-cross-anomaly}, for any \(|S|\ge 2\) and any \(\chi_S\in\prod_{i\in S}\widehat{G_i}\), if there exists some \(i\in S\) such that \(\chi_i=\chi_{i,0}\) is the trivial character, then
\[
\boxed{\;
\Delta_S(\chi_S;\theta)=0\; }.
\]
\end{proposition}
\begin{proof}
If \(\chi_i\) is trivial at some coordinate, then for any \(T\subseteq S\) in the alternating sum of Definition \ref{def:pom-higher-cross-anomaly},
the terms \(\log\mathfrak{M}_{\chi^{(T)}}(\theta)\) and \(\log\mathfrak{M}_{\chi^{(T\cup\{i\})}}(\theta)\) coincide (since the \(i\)-th coordinate takes the trivial character in both cases), while their signs in the sum are opposite, so they cancel pairwise.
Iterating over all pairs that include/exclude \(i\) yields \(\Delta_S(\chi_S;\theta)=0\). \qed
\end{proof}

\begin{theorem}[Quotient space coordinate representation of the Euler defect: higher-order cross-anomalies as minimal linear invariants]\label{thm:pom-euler-defect-quotient-coordinate}
Fix \(\theta\), let \(\widehat G=\prod_{i=1}^k\widehat{G_i}\), and define the function space
\[
\mathcal{F}:=\{f:\widehat G\to\RR:\ f(\chi^0)=0\},
\qquad
\chi^0:=(\chi_{1,0},\dots,\chi_{k,0}).
\]
Let the ``axis-decomposable'' subspace be
\[
\mathcal{S}:=
\Bigl\{f\in\mathcal{F}:\ \exists\,L_i:\widehat{G_i}\to\RR,\ L_i(\chi_{i,0})=0,\ f(\chi)=\sum_{i=1}^k L_i(\chi_i)\Bigr\}.
\]
For each \(|S|\ge 2\), let \(V_S\) denote the function space on \(\prod_{i\in S}\widehat{G_i}\), subject to the constraint of vanishing whenever any coordinate is trivial (Proposition \ref{prop:pom-cross-anom-support-constraint}).
Define the linear map
\[
\Phi:\mathcal{F}\longrightarrow \bigoplus_{|S|\ge 2}V_S,\qquad
\Phi(f):=\bigl(\Delta_S(\cdot;\theta)\bigr)_{|S|\ge 2},
\]
where \(\Delta_S\) is given by the generalization of Definition \ref{def:pom-higher-cross-anomaly} (applied to \(f=\log\mathfrak{M}_\chi(\theta)\)).
Then:
\begin{enumerate}
  \item \(\ker\Phi=\mathcal{S}\). Equivalently, all higher-order cross-anomalies vanish if and only if \(f\) is axis-decomposable (the linearized formulation of Theorem \ref{thm:pom-axis-decomposability-iff}).
  \item \(\Phi\) induces a quotient space isomorphism
  \[
  \boxed{\;
  \overline\Phi:\mathcal{F}/\mathcal{S}\xrightarrow{\ \sim\ }\bigoplus_{|S|\ge 2}V_S\; }.
  \]
  \item Each coset \([f]\in\mathcal{F}/\mathcal{S}\) admits a unique representative \(f^{\mathrm{int}}\) satisfying that all one-body terms (single-coordinate differences) vanish identically, i.e.,
  \(
  \Delta_{\{i\}}(\cdot;\theta)\equiv 0
  \)
  for all \(i\), and
  \[
  f^{\mathrm{int}}(\chi)=\sum_{|S|\ge 2}\Delta_S(\chi_S;\theta).
  \]
\end{enumerate}
\end{theorem}
\begin{proof}
(1) By Theorem \ref{thm:pom-axis-decomposability-iff}, axis-decomposability holds if and only if \(\Delta_S\equiv 0\) for all \(|S|\ge 2\), which is precisely \(\ker\Phi=\mathcal{S}\).
\par
(2) By (1), \(\Phi\) induces an injection on the quotient space.
Surjectivity follows from an explicit construction: given any family \((g_S)_{|S|\ge 2}\) (where \(g_S\in V_S\)), set
\(
f(\chi):=\sum_{|S|\ge 2} g_S(\chi_S)
\).
By the difference expression in Definition \ref{def:pom-higher-cross-anomaly} (following the proof strategy of Theorem \ref{thm:pom-mobius-complete-decomposition}), when computing \(\Delta_S\) for any \(|S|\ge 2\), only terms contained in \(S\) can contribute; moreover, since \(g_T\in V_T\) vanishes when any coordinate is trivial, the \(T\neq S\) terms cancel under the \(S\)-difference, ultimately yielding \(\Delta_S=g_S\).
\par
(3) For any \(f\in\mathcal{F}\), Theorem \ref{thm:pom-mobius-complete-decomposition} gives
\(
f=\sum_{i}\Delta_{\{i\}}+\sum_{|S|\ge 2}\Delta_S
\).
Setting
\(
f^{\mathrm{int}}:=f-\sum_i \Delta_{\{i\}}
\),
we have \(f^{\mathrm{int}}\equiv f\ (\mathrm{mod}\ \mathcal{S})\) with all one-body terms vanishing, and the stated representation holds.
If another representative \(f'\) also satisfies \(\Delta_{\{i\}}\equiv 0\), then \(f-f'\in\mathcal{S}\) with all single-coordinate components equal to zero, forcing \(f-f'=0\); hence uniqueness holds. \qed
\end{proof}

\begin{corollary}[Dimension formula (finite abelian axes)]\label{cor:pom-euler-defect-dimension-formula}
Under the setting of Theorem \ref{thm:pom-euler-defect-quotient-coordinate}, if each axis is a finite abelian group with
\(
n_i:=|\widehat{G_i}|=|G_i|
\), then
\[
\boxed{\;
\dim(\mathcal{F}/\mathcal{S})
\;=\;
\prod_{i=1}^k n_i - 1 - \sum_{i=1}^k (n_i-1)
\;=\;
\sum_{|S|\ge 2}\prod_{i\in S}(n_i-1)\; }.
\]
\end{corollary}
\begin{proof}
\(\dim\mathcal{F}=\prod_i n_i-1\) (the normalization condition removes one degree of freedom).
Meanwhile, \(\mathcal{S}\cong\bigoplus_{i=1}^k\{L_i:\widehat{G_i}\to\RR,\ L_i(\chi_{i,0})=0\}\), so
\(\dim\mathcal{S}=\sum_i(n_i-1)\).
The final identity follows from
\(
\prod_i(1+(n_i-1))=\sum_{S\subseteq[k]}\prod_{i\in S}(n_i-1)
\)
after removing the \(S=\varnothing\) and singleton terms. \qed
\end{proof}

\begin{theorem}[Orthogonal Decomposition of the Euler Defect: Pythagoras Identity for Cross-Term Energy]\label{thm:pom-euler-defect-orthogonal-pythagoras}
Under the setting of Theorem \ref{thm:pom-euler-defect-quotient-coordinate}, endow \(\widehat G=\prod_{i=1}^k\widehat{G_i}\) with the uniform measure, and view \(\mathcal F\subset L^2(\widehat G)\) as a Hilbert subspace.
For any \(f\in\mathcal F\), define the coordinate averaging operator
\[
(E_i f)(\chi):=\mathbb{E}_{\xi_i\sim \mathrm{Unif}(\widehat{G_i})} f(\chi_1,\dots,\chi_{i-1},\xi_i,\chi_{i+1},\dots,\chi_k),
\qquad
D_i:=I-E_i.
\]
For any subset \(S\subseteq[k]\), define
\(
D_S:=\prod_{i\in S}D_i
\),
\(
E_{S^c}:=\prod_{i\notin S}E_i
\), and set
\[
f_S:=E_{S^c}D_S f.
\]
Then:
\begin{enumerate}
  \item There exists a strict orthogonal decomposition
  \[
  \boxed{\;
  f=\sum_{S\subseteq[k]} f_S,\qquad
  \langle f_S,f_T\rangle_{L^2}=0\ (S\ne T)\;}
  \]
  satisfying the Pythagoras identity
  \[
  \|f\|_{L^2}^2=\sum_{S\subseteq[k]}\|f_S\|_{L^2}^2.
  \]
  \item For each \(\emptyset\ne S\subseteq[k]\) and each \(i\in S\), the zero-mean condition \(E_i f_S=0\) holds; moreover, \(f_{\varnothing}=E_{[k]}f\) is a constant function (equal to the global mean of \(f\)).
\end{enumerate}
\end{theorem}
\begin{proof}
Since \(E_i\) is the \(L^2\) conditional expectation projection onto the subspace ``independent of the \(i\)-th axis,'' the \(E_i\) operators pairwise commute and are self-adjoint idempotents; hence \(D_i=I-E_i\) also pairwise commutes and is self-adjoint.
Expanding the identity
\[
I=\prod_{i=1}^k(E_i+D_i)=\sum_{S\subseteq[k]}E_{S^c}D_S
\]
and applying to \(f\) yields the decomposition formula.
If \(S\ne T\), take some \(i\in S\triangle T\); then \(f_S\) can be written as the image of \(D_i(\cdot)\), while \(f_T\) as the image of \(E_i(\cdot)\); since \(E_i\) is an orthogonal projection (equivalently, \(\langle D_i g, E_i h\rangle_{L^2}=0\)), we obtain \(\langle f_S,f_T\rangle_{L^2}=0\), yielding orthogonality and the Pythagoras identity.
For \(i\in S\), by definition \(f_S\) contains the factor \(D_i\), so \(E_i f_S=E_iD_i(\cdot)=0\). \qed
\end{proof}

\begin{theorem}[Minimal Action Gauge: Cross-Term Energy Equals Minimal Distance to Additive Subspace]\label{thm:pom-euler-defect-min-action-gauge}
Following the notation of Theorem \ref{thm:pom-euler-defect-orthogonal-pythagoras}, define the additive (axis-decomposable) subspace
\[
\mathcal S_{\mathrm{add}}
:=
\Bigl\{g\in L^2(\widehat G):\ \exists\,c\in\RR,\ \exists\,g_i:\widehat{G_i}\to\RR,\ g_i(\chi_{i,0})=0,\ g(\chi)=c+\sum_{i=1}^k g_i(\chi_i)\Bigr\}.
\]
Then the \(L^2\) orthogonal projection onto \(\mathcal S_{\mathrm{add}}\) is
\[
\boxed{\ \Pi_{\mathcal S_{\mathrm{add}}} f=f_{\varnothing}+\sum_{|S|=1} f_S\ }.
\]
Moreover, the minimal residual energy satisfies
\[
\boxed{\ \inf_{g\in\mathcal S_{\mathrm{add}}}\|f-g\|_{L^2}^2=\sum_{|S|\ge 2}\|f_S\|_{L^2}^2\ }.
\]
The right-hand side therefore provides a canonical ``Euler defect energy'': it vanishes if and only if \(f\in\mathcal S_{\mathrm{add}}\) (no cross-terms).
Furthermore, the axis-decomposable subspace of Theorem \ref{thm:pom-euler-defect-quotient-coordinate} satisfies \(\mathcal S=\mathcal S_{\mathrm{add}}\cap\mathcal F\).
\end{theorem}
\begin{proof}
By Theorem \ref{thm:pom-euler-defect-orthogonal-pythagoras}, the Hoeffding components \(\{f_S\}_{S\subseteq[k]}\) of \(f\) are pairwise orthogonal.
Denote
\[
\mathcal H_{\le 1}
:=\mathrm{Im}(E_{[k]})\ \oplus\ \bigoplus_{i=1}^k \mathrm{Im}(E_{[k]\setminus\{i\}}D_i),
\]
where \(\mathrm{Im}(E_{[k]})\) is the constant function subspace.
By definition, \(\mathcal H_{\le 1}\subseteq \mathcal S_{\mathrm{add}}\), and for any \(g\in\mathcal S_{\mathrm{add}}\), its second-order and higher cross-components must vanish (i.e., for all \(|S|\ge 2\), \(g_S=0\)), so \(g=\sum_{|S|\le 1} g_S\in \mathcal H_{\le 1}\).
Therefore, \(\mathcal S_{\mathrm{add}}=\mathcal H_{\le 1}\).
Hence \(f_{\varnothing}+\sum_{|S|=1}f_S\) is precisely the orthogonal projection of \(f\) onto the closed subspace \(\mathcal S_{\mathrm{add}}\), and the minimal distance squared is given by the orthogonal complement energy sum, namely
\(
\sum_{|S|\ge 2}\|f_S\|_{L^2}^2
\).
Finally, ``energy is zero'' if and only if all \(|S|\ge 2\) components vanish, i.e., \(f\in\mathcal S_{\mathrm{add}}\). \qed
\end{proof}
